% P44 Figure: GHET five descriptions hub
\begin{figure}[htbp]
\centering
\begin{tikzpicture}[scale=1.12, every node/.style={font=\small}]
  \node[ugpnode,fill=gray!15,minimum width=2.8cm,minimum height=1.15cm] (hub) at (0,0)
    {\textbf{GTE encoding}\\[-2pt]\scriptsize (equivalent)};
  % Asymmetric radii: NE/NW pushed farther (wide central box)
  \node[ugpnode,fill=blue!8,minimum width=2.3cm] (t1) at (0,4.8) {$T_1$\\[-2pt]\scriptsize$\mathcal{L}[\Phi;g]$};
  \node[ugpnode,fill=blue!8,minimum width=2.3cm] (t2) at (5.8,2.0) {$T_2$\\[-2pt]\scriptsize RS $[5,3,3]_7$};
  \node[ugpnode,fill=blue!8,minimum width=2.3cm] (t3) at (3.2,-3.9) {$T_3$\\[-2pt]\scriptsize RT entropy};
  \node[ugpnode,fill=blue!8,minimum width=2.3cm] (t4) at (-3.2,-3.9) {$T_4$\\[-2pt]\scriptsize MDL min.};
  \node[ugpnode,fill=blue!8,minimum width=2.3cm] (t5) at (-5.8,2.0) {$T_5$\\[-2pt]\scriptsize QEC code};
  \foreach \t in {t1,t2,t3,t4,t5} {
    \draw[ugparr] (hub) -- (\t);
  }
  \draw[gray!60,thin] (t1) -- (t2) -- (t3) -- (t4) -- (t5) -- (t1);
  \node[font=\normalsize\itshape,align=center] at (0,-6.05)
    {twenty directed implications $T_i\Rightarrow T_j$;\\[2pt]
     $T_1\!\Rightarrow\!T_2$ and $T_2\!\Leftrightarrow\!T_5$ Lean-certified};
\end{tikzpicture}
\caption{Five a priori distinct descriptions of the GTE holographic encoding
  structure are pairwise equivalent (Theorem~\ref{thm:ghet}). Any single
  description determines the full content.}
\label{fig:ghet_five_descriptions}
\end{figure}
